\section{Gamma coefficients prove every harmonic-order identity}
\label{ghharm:sec:main}
The continuous Gamma correlations above have a discrete counterpart.
Finite harmonic products are coefficients of a Gamma ratio; a normally
convergent Gauss sum therefore evaluates an entire family at once.
This gives corrected all-order forms of the fifth- and sixth-order
conjectures in Choi~\cite{ghharm:Choi}, rather than isolated numerical
relations. The harmonic-resonance continuation supplies this route and
its broader mixed-jet context~\cite{ghharm:report}.

For $n,r\ge0$, define
\begin{equation}\label{ghharm:eq:finite-product}
 \prod_{j=1}^n\left(1+\frac uj\right)
   =\sum_{r=0}^n\frac{\mathcal E_r(n)}{r!}u^r,
 \qquad \mathcal E_0(n)=1.
\end{equation}
Thus $\mathcal E_r(n)$ is $r!$ times the elementary symmetric
polynomial in $1,1/2,\ldots,1/n$. In particular
$\mathcal E_1(n)=H_n$ and
$\mathcal E_2(n)=H_n^2-H_n^{(2)}$. Equivalently it is the complete
exponential Bell polynomial with arguments
$H_n,-H_n^{(2)},2H_n^{(3)},\ldots,
(-1)^{r-1}(r-1)!H_n^{(r)}$.

\begin{theorem}[Every pair of harmonic coefficient orders]
\label{ghharm:thm:generator}
For $u,v$ in a neighborhood of zero,
\begin{equation}\label{ghharm:eq:generator}
 \sum_{n=0}^{\infty}
 \frac{(1+u)_n(1+v)_n}{(n!)^2(n+1)(n+2)}
 =\frac{\Gamma(1-u-v)}{\Gamma(2-u)\Gamma(2-v)}.
\end{equation}
Every fixed Taylor coefficient can be taken termwise. Consequently,
for all $r,t\ge0$,
\begin{align}\label{ghharm:eq:mixed}
 \sum_{n=1}^{\infty}
 \frac{\mathcal E_r(n)\mathcal E_t(n)}{(n+1)(n+2)}
 ={}&r!t![u^rv^t]
       \frac{\Gamma(1-u-v)}{\Gamma(2-u)\Gamma(2-v)}
       -\frac12\mathbf1_{r=t=0}.
\end{align}
These are ordinary absolutely convergent series.
\end{theorem}
\begin{proof}
Since $(3)_n=n!(n+1)(n+2)/2$, the series in
\eqref{ghharm:eq:generator} is
$\tfrac12{}_2F_1(1+u,1+v;3;1)$. Gauss summation
\cite{ghharm:Gauss} evaluates it as the displayed Gamma quotient.
On a sufficiently small closed polydisc, its summands are uniformly
$O(n^{-2+\operatorname{Re}(u+v)})$ by the Gamma-ratio asymptotic.
Enlarging that polydisc slightly and applying Cauchy's coefficient
formula justifies each derivative or coefficient extraction.
Formula~\eqref{ghharm:eq:finite-product} now proves
\eqref{ghharm:eq:mixed}. The $n=0$ summand is $1/2$ and contributes
only at $r=t=0$.
\end{proof}

\begin{corollary}[Two explicit all-order families]
\label{ghharm:cor:two-families}
For every $r\ge1$ and every $t\ge0$, respectively,
\begin{align}
 \sum_{n=1}^{\infty}\frac{\mathcal E_r(n)}{(n+1)(n+2)}
   &=r!,\label{ghharm:eq:pure}\\
 \sum_{n=1}^{\infty}\frac{H_n\mathcal E_t(n)}{(n+1)(n+2)}
   &=t!\left(1+\sum_{j=2}^{t+1}\zeta(j)\right).
   \label{ghharm:eq:one-harmonic}
\end{align}
The sum on the second right side is empty at $t=0$.
\end{corollary}
\begin{proof}
Set $v=0$ in \eqref{ghharm:eq:generator}. Its right side is
$1/(1-u)$, proving \eqref{ghharm:eq:pure}.
Differentiate first in $v$ and then put $v=0$. The result is
\[
 \frac{\psi(2)-\psi(1-u)}{1-u}
   =\frac{1+\sum_{j\ge1}\zeta(j+1)u^j}{1-u},\qquad |u|<1.
\]
The normally convergent digamma Taylor expansion at one proves this
last equality. Taking the coefficient of $u^t$ proves
\eqref{ghharm:eq:one-harmonic}. Its $n=0$ derivative is zero, so no
additional constant occurs.
\end{proof}

The cases $r=5,6$ prove the intended pure harmonic identities in
Choi's Conjecture~4; the cases $t=5,6$ prove corrected versions of
Conjecture~6. His printed formulas (41)--(46) carry twice the
multiplier in \eqref{ghharm:eq:one-harmonic}. The finite-product
definition also fixes the harmonic orders in every numerator, including
the sixth-order term $-120H_n^{(6)}$ and, after multiplication by
$H_n$, the partition term $45H_n^3(H_n^{(2)})^2$.
The unambiguous coefficient formula is the corrected statement at all
orders. The source comparison and its text-versus-visual evidence
boundary are recorded in \src{verification/ninth-analytic-audit.md}.

For a short independent check of the simplest normalization, summation
by parts gives
\begin{align*}
 \sum_{n\ge1}\frac{H_n^2}{(n+1)(n+2)}
 &=\sum_{k\ge1}\frac{H_k^2-H_{k-1}^2}{k+1}\\
 &=2\sum_{k\ge1}\frac{H_{k-1}}{k(k+1)}
       +\sum_{k\ge1}\frac1{k^2(k+1)}\\
 &=2+\zeta(2)-1=1+\zeta(2).
\end{align*}
Indeed Tonelli's theorem gives
$\sum_{k\ge1}H_{k-1}/(k(k+1))
=\sum_{j\ge1}1/(j(j+1))=1$.
This isolates the factor of two without any special-function
differentiation. Higher mixed orders are finite Gamma Taylor
coefficients; no arithmetic-independence assumption is needed.
